% =====================================================================
% SST Canon v0.8.35 candidate addendum
% Seven-article closure, attribution, and inference guards
% Status: main-canon methodology/orthodox identities only
% =====================================================================

\subsection{Closure, Attribution, and Representation Guards}
\label{subsec:seven_article_closure_guards}

\paragraph{Status and scope.}
\textbf{[MIXED STATUS / STATEMENT-LEVEL TAGS BELOW].}
This subsection imports no acoustic, optical, magnetic-vortex, or
Ho\v{r}ava--Lifshitz model as SST dynamics.  It records orthodox continuum
identities and inference rules that any SST realization must respect.

\paragraph{General incompressible pressure closure.}
\textbf{[ORTHODOX / EXACT BULK IDENTITY].}
For a constant-density incompressible velocity field
\(\partial_i u_i=0\), define
\begin{align}
A_{ij}&:=\partial_j u_i,\\
S_{ij}&:=\frac12(A_{ij}+A_{ji}),\\
\boldsymbol\omega&:=\nabla\times\mathbf u.
\end{align}
Taking the divergence of the Euler momentum equation gives
\begin{align}
\boxed{
\nabla^2\!\left(\frac{p}{\rho_{\!f}}\right)
=\frac12\lVert\boldsymbol\omega\rVert^2-S_{ij}S_{ij}
}
\label{eq:canon_pressure_poisson_enstrophy_strain}
\end{align}
throughout a smooth bulk region.  Consequently a local vorticity maximum
is not by itself sufficient to identify a pressure minimum: the strain
invariant and the nonlocal boundary-value problem are part of the same
closure.

For a free-space decaying solution the corresponding Green representation is
\begin{align}
\boxed{
\frac{p(\mathbf x)-p_\infty}{\rho_{\!f}}
=-\frac{1}{4\pi}\int_{\mathbb R^3}
\frac{\frac12\lVert\boldsymbol\omega(\mathbf x')\rVert^2-S:S(\mathbf x')}
{\lVert\mathbf x-\mathbf x'\rVert}\,d^3x'
}
\label{eq:canon_pressure_green_enstrophy_strain}
\end{align}
when the stated decay and regularity conditions hold.  Bounded or periodic
problems require the Green function appropriate to their boundary conditions.

\paragraph{Bernoulli scope guard.}
\textbf{[ORTHODOX / CRITICAL NOTE].}
A relation such as
\(\delta p=-\rho_{\!f}\delta|\mathbf u|^2/2\) is retained only under the
stationary Bernoulli assumptions declared where it is used.  It is not a
general three-dimensional pressure closure for an evolving knotted vortex.
Equation~\eqref{eq:canon_pressure_poisson_enstrophy_strain} is the bulk
constraint that must close any general incompressible reconstruction.

\paragraph{Topology-layer non-equivalence.}
\textbf{[CANON INFERENCE GUARD].}
The following objects are distinct unless an explicit theorem or resolved
field calculation proves an identification:
\begin{align}
\mathcal K_{\rm center},\qquad
\mathcal K_{\omega},\qquad
\mathcal K_{\Phi},\qquad
\mathcal K_{\rm null},
\end{align}
where the labels denote centerline, material-vorticity-tube, internal-phase
fiber, and observable-null/singularity topology respectively.  A central
phase or amplitude null is therefore not sufficient evidence for a material
threaded vortex core, and a phase-fiber knot does not by itself classify the
centerline knot.

\paragraph{Effective-metric derivation guard.}
\textbf{[ORTHODOX METHOD / SST THEOREM TARGET].}
An SST effective metric for a probe may be introduced only after the
linearized probe equations have supplied a hyperbolic principal symbol
\begin{align}
\mathcal P(x,k)=g_{\rm eff}^{\mu\nu}(x)k_\mu k_\nu
\end{align}
(or an explicitly stated anisotropic generalization).  The resulting
\(g_{\rm eff}^{\mu\nu}\) is first a probe-response geometry; it is not by
itself a derivation of curvature of the underlying Euclidean substrate or of
the Einstein field equations.

\paragraph{Representation-invariance rule.}
\textbf{[CANON METHODOLOGY].}
If two formulations are claimed to be exactly equivalent descriptions of
the same on-shell SST state, every physical observable must agree after the
correct variable transformation.  A closed-loop phase, frequency, or
orientation defect produced by changing representations is not promoted as
physical holonomy unless it survives numerical refinement and cannot be
accounted for by the incompressibility residual, Euler residual, discretization,
or an omitted transformation term.

\paragraph{Kinematics-versus-dynamics rule.}
\textbf{[CANON INFERENCE GUARD].}
A fitted kinematic ansatz---including a clock reparameterization, dispersion
curve, cosmographic expansion history, or empirical phase warp---is not
evidence for the underlying SST dynamics unless the mapping from SST
primitives to the fit parameters is derived independently.  Derived
observables must be reproducible from the frozen solution/posterior, and the
model must pass domain-regularity, asymptotic-limit, and complexity-penalized
comparison tests before a dynamical interpretation is promoted.
